% Part: proof-theory
% Chapter: proof-search
% Section: introduction

\documentclass[../../../include/open-logic-section]{subfiles}

\begin{document}

\olfileid{pt}{ps}{int}

\olsection{Pambuka}

Aksioma lan aturan inferènsi ing sistem !!{proof} kayata
\Log{G3c} utawa \Log{N1i} nemtokaké wit sekuèn utawa
!!{formula}s kang dianggep !!{proof}s bener. Yèn wis
diwènèhi wit sekuèn utawa !!{formula}s kaya mangkono
(bokmanawa kanthi katerangan tambahan ngenani aturan kang
ditrapaké ing saben langkah, utawa label kanggo asumsi
lan aturan inferènsi kang mbusak asumsi iku), définisi
aksioma lan aturan ngidini kita mriksa apa wit iku
!!{proof} kang bener. Nanging, \emph{nemokaké}
!!a{proof} tumrap sekuèn tartamtu, utawa tumrap
!!{formula} tartamtu saka asumsi kang diwènèhaké,
iku masalah liya kang lumrahé luwih angel. Masalah
iki diarani \emph{panelusuran bukti}.

Ana tata cara kang prasaja banget kanggo panelusuran bukti.
!!{formula}s bisa dianggep rerangkéyan lambang saka
alfabèt winates. Sekuèn lan wit sekuèn utawa !!{formula}s
uga bisa dianggep rerangkéyan !!{formula}s (bokmanawa
nganggo kurung mligi kanggo nandhani percabangan wit).
Aksioma lan aturan ngidini kita mriksa kanthi mekanis
apa rerangkéyan lambang minangka !!{proof} kang bener.
Mula kita bisa ngasilaké kanthi mekanis \emph{kabèh
!!{proof}s kang mungkin lan bener}: ngasilaké kabèh
rerangkéyan lambang saka alfabèt kang dipigunakaké
kanggo makili !!{proof}s, banjur mriksa saben
rerangkéyan calon apa makili !!{proof} kang bener.
Tata cara prasaja iki ngasilaké kabèh !!{proof}s bener
nganti nemokaké siji kang dadi !!{proof} tumrap sekuèn
kang diwènèhaké (utawa tumrap !!{formula} kang asumsi
mbukaké kalebu ing asumsi kang diwènèhaké).

Cara kang luwih becik yaiku cara lumrah nalika nggoleki
!!a{proof} kanthi tangan: miwiti saka sekuèn pungkasan,
milih !!a{formula}, banjur ngetrapaké aturan inferènsi
``mundur'' kanggo ngasilaké premis siji utawa sawetara.
Yèn premis-premis iku nduwèni !!{proof}s, premis mau
bisa digabung karo inferènsi pungkasan dadi
!!a{proof} tumrap sekuèn kang diwènèhaké. Cara kang
padha, senajan luwih ruwet, bisa dienggo nggoleki
!!a{proof} ing sistem deduksi alamiah.

Nanging cara iki ora mesthi kasil: yèn milih aturan kang
ora trep utawa urutan !!{formula}s kang ora trep, kita
bisa tekan dalan buntu. Carané uga durung ditemtokaké
kanthi pepak: ing saben tataran bisa ana luwih saka siji
!!{formula} kang bisa dipilih, utawa luwih saka siji
aturan kang bisa ditrapaké mundur. \emph{Algoritma
panelusuran bukti} yaiku tata cara kang wis ditemtokaké
kanthi pepak lan bakal nemokaké !!a{proof} yèn pancèn
ana (yaiku ora bakal klèru amarga pilihan langkah).

Sawetara sistem !!{proof} luwih gampang digarap nganggo
algoritma panelusuran bukti tinimbang sistem liyané.
Ing sistem sekuèn, !!{main operator} saka
!!a{formula}, lan panggonané ing sisih kiwa utawa
tengen, nemtokaké aturan kang kudu ditrapaké mundur.
Tuladhané, yèn arep nemokaké !!a{proof} tumrap
$\Gamma \Sequent \Delta, !A \land !B$ kanthi
$!A \land !B$ ing sisih tengen, kita ngetrapaké
\RightR{\land} mundur lan ngasilaké loro premis kang
cocog: $\Gamma \Sequent \Delta, !A$ lan
$\Gamma \Sequent \Delta, !B$. Nanging, ing sistem
\Log{G1c}, anané kontraksi gawé prakara iki luwih
angel amarga pilihané luwih akèh: bisa uga ngetrapaké
\RightR{\Contraction} mundur kanggo ngasilaké premis
$\Gamma \Sequent \Delta, !A \land !B, !A \land !B$.
Sabanjuré bisa ngasilaké premis
$\Gamma \Sequent \Delta, !A \land !B, !A \land !B, !A \land !B$,
lan saterusé. Ing \Log{G2c} luwih angel manèh.
Ngetrapaké \RightR{\land} mundur uga mbutuhaké tebakan
$\Gamma_1$, $\Gamma_2$ saéngga
$\Gamma = \Gamma_1 \cup \Gamma_2$, lan
$\Delta_1$, $\Delta_2$ saéngga
$\Delta = \Delta_1 \cup \Delta_2$; banjur nyoba premis
$\Gamma_1 \Sequent \Delta_1, !A$ lan
$\Gamma_2 \Sequent \Delta_2, !B$. Tebakan kang salah
bisa nuwuhaké dalan buntu mengko, mula kudu bali menyang
wiwitan lan milih $\Gamma_1$, $\Gamma_2$,
$\Delta_1$, $\Delta_2$ kang béda. Yèn !!{formula}
iku $!A \lor !B$, kita kudu milih salah siji saka
loro kemungkinan \RightR{\lor}, kang ngasilaké
salah siji premis $\Gamma \Sequent \Delta, !A$ utawa
$\Gamma \Sequent \Delta, !B$. Pilihan kang salah
uga bakal meksa kita bali menyang langkah sadurungé.

Sistem \Log{G3c} ora nduwèni masalah iki. Sistem iku
ora nduwèni aturan struktural, lan aturan inferènsi
ditetepaké déning !!{main operator} lan sisih
panggonané !!{formula}. Mula \Log{G3c} luwih cocog
kanggo panelusuran bukti tinimbang \Log{G1c} utawa
\Log{G2c}. Panelusuran kang kasil bakal ngasilaké
!!a{proof} ing \Log{G3c}; !!{proof} iku banjur bisa
diterjemahaké dadi bukti ing \Log{G1c} utawa
\Log{G2c} (utawa malah \Log{N1c}). Mula, algoritma
panelusuran bukti kanggo \Log{G3c}, bebarengan karo
algoritma panerjemah \Log{G3c}-!!{proof}s dadi
!!{proof}s ing sistem liya, uga ngrampungaké
panelusuran bukti kanggo sistem-sistem mau.

Kepiyé carané ngerti yèn algoritma panelusuran bukti
ora bakal gagal nalika bukti pancèn ana? Kita kudu
nuduhaké yèn nalika algoritma ora nemokaké
!!a{proof}, ora ana !!{proof} kang bisa ana.
Algoritma kang ora becik bisa gagal nemokaké bukti
senajan buktiné ana. Tata cara prasaja mau ora nduwèni
masalah iki, amarga mriksa kabèh !!{proof}s kang mungkin.
Nanging algoritma panelusuran bukti diarepaké luwih
éfisièn lan mung ngentèkaké gawé sakperluné.

Gagasan pokok kanggo nuduhaké kasilé algoritma
panelusuran bukti yaiku: yèn algoritma gagal, saka
carané gagal kita bisa nuduhaké yèn sekuèn mau ora
mungkin nduwèni !!a{proof}. Carané yaiku mbuktèkaké
yèn sekuèn iku ora valid. Ing logika orde kapisan
klasik, sekuèn $\Gamma \Sequent \Delta$ valid yèn
saben !!{structure} ndadèkaké paling ora siji
!!{formula} ing $\Gamma$ salah utawa paling ora
siji !!{formula} ing~$\Delta$ bener. \Log{G3c}
\emph{sahih}, yaiku mung !!{prove}s sekuèn kang valid.
Iki dibuktèkaké kanthi induksi miturut !!{proof}s:
aksioma cetha valid, lan yèn premis aturan valid,
konklusiné uga valid. Kanggo nuduhaké yèn algoritma
panelusuran bukti mesthi kasil nalika bukti ana,
kita nuduhaké yèn panelusuran !!a{proof} tumrap
$\Gamma \Sequent \Delta$ gagal, kita bisa netepaké
!!a{structure} kang ndadèkaké kabèh !!{formula}s
ing $\Gamma$ bener lan kabèh !!{formula}s ing~$\Delta$
salah. Iki uga netepaké yèn \Log{G3c}
\emph{lengkap}: yèn $\Gamma \Sequent \Delta$ valid,
ana \Log{G3c}-!!{proof} tumrap sekuèn iku.
Amarga saben panelusuran gagal nuduhaké yèn
$\Gamma \Sequent \Delta$ ora valid, saben panelusuran
tumrap sekuèn valid mesthi kasil.

\end{document}
